\section{Excluding nonabelian simple colors}\label{sec:palettes}

We retain the odd-degree division-algebra hypotheses and notation of
Section~\ref{sec:colors}. Thus $L/k$ is quadratic, $D$ is a central division
$L$-algebra of odd degree $n\geq3$, the involution $*$ induces the nontrivial
automorphism of $L/k$, and $S=\SU(D,*)\subset U=\U(D,*)$.
Fix $i\in L^\times$ with $\bar i=-i$.
For $R=S(k)^e$ let $P_0$ and $V_0$ have the meaning of
Proposition~\ref{prop:finite-defect}. We now eliminate the nonabelian
alternative in Proposition~\ref{prop:color-dichotomy}.

\begin{theorem}\label{thm:no-simple-colors}
There is no homomorphism
\[
 \phi:U(k)\longrightarrow\mathcal F
 \quad\text{with}\quad
 R\subset\ker\phi,\qquad \phi(V_0)=\mathcal F,
\]
where $\mathcal F$ is a finite nonabelian simple group.
\end{theorem}

Suppose, for a contradiction, that such a homomorphism exists.
We call its values \emph{colors}. It is surjective, so central scalars in
$U(k)$ have color $1$: their images centralize $\mathcal F$.
The finite-group input is the following assertion. Its statement is
independent of the arithmetic constructions below.

\begin{lemma}\label{lem:finite-obstruction}
For every finite nonabelian simple group $\mathcal F$ there is a nonempty
proper conjugacy-invariant subset $\mathcal S\subset\mathcal F$ with the
following property. If $W\in\mathcal S$, $Z\notin\mathcal S$, and $WZ=ZW$,
then it is impossible that
\begin{equation}\label{eq:palette-obstruction}
 ZW\in C(B)(ZW^{-1})C(BZ^2)
 \qquad\text{for every }B\in\mathcal F,
\end{equation}
where $C(B)=C_{\mathcal F}(B)$ denotes the centralizer of $B$ in
$\mathcal F$.
\end{lemma}

The proof of Lemma~\ref{lem:finite-obstruction} occupies
Section~\ref{sec:finite-groups}. We use exactly its displayed
double-coset condition. Our task in this section is to force that condition
from sufficiently many additive returns of colors, and then to show that
avoiding it would make membership in $\mathcal S$ constant on
$\mathcal F$.

\subsection{The compact space of color configurations}

Let $\widetilde G$ be the simply connected special unitary group of the
hyperbolic binary hermitian form over $(D,*)$ with coefficients $1,-1$.
Let $X$ be its projective variety of isotropic right $D$-lines.
Since $D$ is division, the form has relative rank one, and the stabilizer
of such a line is a minimal $k$-parabolic of $\widetilde G$.
Every rational isotropic line has a unique representative $(u,1)$ with
$u\in U(k)$: the second entry cannot vanish, and division makes it
invertible. This identifies $X(k)$ with $U(k)$.

There is also an additive parametrization. Put
\[
 J=\{x\in D:x=x^*\},\qquad
 q(x)=(x+i)(x-i)^{-1}.
\]
Here $J$ is a $k$-vector space, and
\begin{equation}\label{eq:palette-cayley}
 X(k)=J(k)\sqcup\{\infty\},\qquad q(\infty)=1.
\end{equation}
Indeed $x-i$ is nonzero for $x=x^*$, and $u-1$ is invertible for every
$u\ne1$ in $U(k)$. The corresponding column on the additive chart is
$(x+i,x-i)$. At a completion $k_v$, these two parametrizations give open
charts isomorphic to $U(k_v)$ and $J(k_v)$, respectively. Their
complements are proper algebraic subvarieties and therefore have measure
zero for any smooth measure on $X(k_v)$. We shall use this assertion at
finitely many completions at a time.

Let $\mathcal H$ be the countable group of rational transformations of
$X$ induced by rational form similitudes. It contains $\widetilde G(k)$,
the left and right rotations
\[
 L_a(u)=au,\qquad R_a(u)=ua\quad(a\in U(k)),
\]
and the additive translations $T_b(x)=x+b$ for $b\in J(k)$.
Write $\tau_s=T_s$ when $s\in k\subset J(k)$.
For $p\in X(k)$ define a configuration
\[
 c(p)=(c_h(p))_{h\in\mathcal H},\qquad
 c_h(p)=\phi(u(hp))\in\mathcal F.
\]
We call this configuration a \emph{palette}, and let
$\mathcal C\subset\mathcal F^{\mathcal H}$ be the closure of the rational
palettes. The product topology makes $\mathcal C$ compact and metrizable.
The action is
\[
 (h c)_j=c_{jh};\qquad c(hp)=h c(p).
\]
Every pattern on finitely many coordinates of a palette in $\mathcal C$
occurs at a rational point, because the colors form a finite discrete
space. In particular, all finite-coordinate identities satisfied by
rational palettes hold throughout $\mathcal C$.

\subsection{A translation-invariant law with uniform marginals}

The next proposition supplies additive recurrence without losing any
color. Arbitrary additive averages of rational palettes need not have
uniform marginals, so we first construct a measure with controlled
conditional laws over the local flag varieties.

\begin{proposition}\label{prop:palette-measure}
There is a probability measure $\mu$ on $\mathcal C$ invariant under
$T_b$ for every $b\in J(k)$ such that
\[
 \mu\{c:c_h=a\}=\frac1{|\mathcal F|}
 \qquad(h\in\mathcal H,\ a\in\mathcal F).
\]
\end{proposition}

\begin{proof}
We construct a conformal joint law, identify its local projection, prove
that its conditional color marginals are uniform, and only then average
additively.

\paragraph{A joint law over all completions.}
For every place $v$, choose an Iwasawa maximal compact subgroup
$K_v\subset\widetilde G(k_v)$ and its invariant probability $m_v$ on
$X(k_v)$. The compact group is transitive on this flag variety.
Choose the good integral compacts at almost all finite places, and set
\[
 Y=\prod_v X(k_v),\qquad m=\prod_v m_v.
\]
Each $m_v$ is smooth and has full support. A fixed $h\in\mathcal H$
has smooth positive Jacobians on all local factors and preserves $m_v$
outside finitely many places. For the latter assertion, discard the
finitely many denominators and nonunit similitude multipliers of a
representative of $h$. At every remaining place it normalizes the good
integral compact of $\widetilde G$; its pushforward of $m_v$ is consequently
the same compact-invariant probability. Thus
\[
 \rho(h,y)=\frac{d(h_*m)}{dm}(y)
\]
is a continuous positive function on $Y$, given by a finite product, and
\begin{equation}\label{eq:palette-cocycle}
 \rho(ag,y)=\rho(a,y)\rho(g,a^{-1}y).
\end{equation}

Assign positive weights $w(p)$ to rational points by requiring the ratio
$w(hp)/w(p)$ to be the product of the forward local volume Jacobians of
$h$ at $p$. To construct them, normalize the weight at one rational
base point and transport by $\widetilde G(k)$. Rational conjugacy of
parabolics of the given type gives transitivity
\cite[Theorem~4.13(a,c)]{BorelTits1965}. If two transporters
give the same point, their quotient fixes the base point; the product
of its tangent determinants is $1$ by the product formula in $k$.
This proves independence of the transporter. The same reasoning at
rational fixed points proves the weight-ratio rule for every
$h\in\mathcal H$, not just for $\widetilde G(k)$.

We shall use
\begin{equation}\label{eq:palette-weight-sum}
 \sum_{p\in X(k)}w(p)^s<\infty\qquad(s>1).
\end{equation}
Here is the convergence theorem and normalization involved. Let $P$ be
the stabilizing minimal $k$-parabolic. For a local decomposition
$g=k_0p_0$ with $k_0\in K_v$ and $p_0\in P(k_v)$, the forward tangent
Jacobian at the base point is $\delta_P(p_0)^{-1}$, where $\delta_P$ is
the modular character of $P$. After replacing $g$ by $g^{-1}$, the
sum in Equation~\ref{eq:palette-weight-sum} is the spherical
minimal-parabolic Eisenstein sum with inducing character $\delta_P^s$.
The derived Levi is $k$-anisotropic, since $P$ is minimal and
$\widetilde G$ has relative rank one. More precisely, its Levi subgroup
$M$ is $k$-anisotropic modulo its maximal $k$-split central torus, so
$M(k)\backslash M(\bbA_k)^1$ is compact, where the superscript $1$
imposes adelic norm one for every $k$-rational character of $M$
\cite[Theorem~5.6]{PRbook}.
The constant inducing function is therefore square-integrable modulo
the split center and cuspidal, since $M$ has no proper $k$-parabolics.
The usual absolute-convergence
theorem applies in the Godement domain
$\operatorname{Re}\lambda>\rho_P$ for normalized parameter
$\rho_P+\lambda$, where $2\rho_P$ is the sum of the positive relative
roots with multiplicities; see
\cite[Chapter~4, Lemma~4.1, and Appendix~II]{Langlands}.
For $\delta_P^s$ this parameter is
$\lambda=(2s-1)\rho_P$, and the domain is exactly $s>1$.
We use the number-field adelic formulation, equivalently the arithmetic
formulation after restriction of scalars and finite adelic double-coset
decomposition.

Give $(p,c(p))\in Y\times\mathcal C$ mass proportional to $w(p)^s$,
and denote this probability by $\mu_s$. Its transformation rule is
\[
 \frac{d(h_*\mu_s)}{d\mu_s}(y,c)=\rho(h,y)^s.
\]
The reciprocal in this formula is worth noting: the mass at $hp$ after
pushforward is the old mass at $p$, so the forward Jacobian is inverted.
Take a weak limit $\mu^0$ as $s\downarrow1$. Compactness gives such a
limit, and uniform convergence of $\rho(h,\cdot)^s$ for each fixed $h$
gives
\begin{equation}\label{eq:palette-joint-conformal}
 h_*\mu^0=\rho(h,\cdot)\mu^0.
\end{equation}

\paragraph{The base projection is the product measure.}
Let $\nu$ be the projection of $\mu^0$ to $Y$. We prove $\nu=m$.
Fix a sufficiently large finite set $T$ of places containing the
archimedean and model-exceptional places. Use subscripts $T$ for products
over $T$, and superscripts $T$ for products over its complement.
An arithmetic lattice $\Gamma\subset\widetilde G_T$, integral outside
$T$, preserves the tail probability. Equation~\ref{eq:palette-joint-conformal}
therefore gives
\[
 \gamma_*\nu_T=\rho_T(\gamma,\cdot)\nu_T
 \qquad(\gamma\in\Gamma).
\]
Define the continuous positive function
\[
 f(g)=\int_{Y_T}\rho_T(g,y)\,d\nu_T(y)
 \qquad(g\in\widetilde G_T).
\]
The cocycle identity and the preceding transformation rule give
$f(\gamma g)=f(g)$.

Choose a full-support probability $\eta$ on $\widetilde G_T$ that is
left invariant by $K_T=\prod_{v\in T}K_v$. Compact transitivity and
$\int\rho_T(h,y)\,dm_T(y)=1$ imply
\[
 \int\rho_T(h,z)\,d\eta(h)=1\quad(z\in Y_T).
\]
Indeed average first over left multiplication of $h$ by $K_T$; the
result averages $\rho_T(h,\cdot)$ over $m_T$.
Equation~\ref{eq:palette-cocycle} now gives
\[
 f(g)=\int f(gh)\,d\eta(h).
\]
This is a positive harmonic function on the finite-volume space
$\Gamma\backslash\widetilde G_T$.
To see it is constant without an unproved integrability assertion, apply
Jensen's inequality to the bounded strictly concave function
$q(t)=t/(1+t)$ for $t>0$. Its nonnegative Jensen discrepancy has integral
zero over the quotient, by right invariance of the quotient volume.
Equality in strict Jensen makes $f(gh)$ constant for $\eta$-almost all
$h$, for almost every coset $g$. Full support of $\eta$ and continuity
then make $f$ constant everywhere. Its value is $f(1)=1$.

We still have to show that these integrals determine $\nu_T$.
Choose a sequence $a_j\in\widetilde G_T$ contracting the big cell of
$Y_T$ to a point $x_0$, using a split cocharacter for the parabolic at
each place. The complement of that cell has $m_T$-measure zero, so
$(a_j)_*m_T\to\delta_{x_0}$. For $b\in C(Y_T)$ put
\[
 B_j(y)=\int_{K_T}b(k_0x_0)\rho_T(k_0a_j,y)\,dk_0.
\]
Under the probability density $\rho_T(k_0a_j,y)\,dk_0$, the variable
$k_0^{-1}y$ has distribution $(a_j)_*m_T$: use compact invariance and
push Haar measure on $K_T$ onto $Y_T$. Compact equicontinuity shows that
$B_j(y)\to b(y)$ uniformly in $y$. On the other hand, $f(k_0a_j)=1$
gives
\[
 \int B_j\,d\nu_T
 =\int_{K_T}b(k_0x_0)\,dk_0
 =\int b\,dm_T.
\]
It follows that $\nu_T=m_T$. Exhausting the places proves $\nu=m$.

Disintegrate the joint law as
\begin{equation}\label{eq:palette-disintegration}
 d\mu^0(y,c)=dQ_y(c)\,dm(y).
\end{equation}
The spaces are compact metrizable, so regular conditional probabilities
$Q_y$ exist. Equation~\ref{eq:palette-joint-conformal} gives
$Q_{hy}=h_*Q_y$ almost everywhere for each $h$. Since $\mathcal H$ is
countable, all these equalities may be imposed on a common conull set.

\paragraph{Uniform conditional color marginals.}
Write $P(y)$ for the probability vector on $\mathcal F$ which is the
$c_1$-marginal of $Q_y$, and put $N=\ker\phi\cap S(k)$.
The exact palette identities
\[
 c_{L_a h}=\phi(a)c_h,\qquad
 c_{R_a h}=c_h\phi(a)
\]
imply that $P$ is invariant under the separate left and right rotations
by every $a\in N$. We first upgrade this invariance at one completion,
using a finite-volume unitary representation, and then use approximation
away from that completion. No topology is placed on the abstract map
$\phi$.

Enlarge $T$ so that it contains a finite place $v_0$ where $S$ has
positive rank. On the $u_T$-chart, $m_T$ is in the Haar measure class
of $U_T$. Disintegrate this class along the determinant map. On a fiber
of determinant $\lambda$ choose a measurable representative
$d_\lambda\in U_T$ and write
\[
 u_T=g d_\lambda,\qquad g\in S_T.
\]
The determinant map is a smooth group map onto its image, so the
conditional measure class along its fibers is Haar measure on $S_T$.
For almost every $\lambda$, the function $P(gd_\lambda,y^T)$ is invariant
under the diagonal left action of
\[
 \Gamma'=N\cap\{a\in S(k):a_v\text{ is integral for }v\notin T\}.
\]
This is a finite-index subgroup of an arithmetic lattice. Consequently
each component of $P$ defines a bounded function on
\begin{equation}\label{eq:palette-lattice-space}
 \Gamma'\backslash(S_T\times Y^T).
\end{equation}
The space has finite invariant measure, obtained from Haar measure on
the group factor and $m^T$ on the tail. For example, a measurable
fundamental domain for $\Gamma'$ in $S_T$, together with the tail,
constructs this measure. Left integrality ensures that the tail action
preserves $m^T$.

Right invariance by $a\in\Gamma'$ acts on
Equation~\ref{eq:palette-lattice-space} by
\[
 (g,y^T)\longmapsto
 (g d_\lambda a_Td_\lambda^{-1},\ R_a y^T).
\]
These actions commute with the diagonal left action. The right tail
rotations lie in a compact group of measure-preserving transformations,
since all their local coordinates are integral outside $T$.

There is a sequence $a_j\in\Gamma'$ escaping every compact set at
$v_0$ and bounded in the other factors of $S_T$. To obtain it, use the
finite-index closure of $N$ and strong approximation off $v_0$ to find
infinitely many elements in a fixed compact cylinder at all the other
places. The full arithmetic diagonal is discrete, so a subsequence
must escape at $v_0$. Pass to a subsequence on which the bounded factors
and the compact tail rotations converge. The fixed conjugations by
$d_\lambda$ preserve these escape and convergence properties.

Let $\mathcal L$ be the $L^2$ space of
Equation~\ref{eq:palette-lattice-space}, and decompose it into the
$S(k_{v_0})$-invariant vectors and their orthogonal complement.
Howe--Moore decay applies on the latter: $S(k_{v_0})$ is an isotropic
simple local group, and local Kneser--Tits identifies it with the group
generated by its root groups; see \cite{HoweMoore} or
\cite[Theorem~4.19 and Definition~4.23]{Ciobotaru}.
For completeness, let $v$ be the projection of a component of $P$ to
this complement. The simultaneous right actions fixing $P$ also fix
$v$. They are products $\pi(g_j)A_j$, where $g_j\to\infty$ in
$S(k_{v_0})$ and the commuting unitary operators $A_j$ converge strongly
to a unitary operator $A$. Therefore
\[
 \|v\|^2
 =\langle\pi(g_j)A_jv,v\rangle
 =\langle\pi(g_j)Av,v\rangle+o(1)\longrightarrow0.
\]
Thus $P$ is invariant under the full right $S(k_{v_0})$ action.
All descent and invariance statements above hold on almost every
determinant fiber by ordinary measure disintegration and countability
of the rational groups. Strong continuity of the local and compact
actions follows first on continuous compactly supported functions and
then on $L^2$.

Full right $S(k_{v_0})$-invariance makes $P(y)$ depend on $u_{v_0}$
only through its determinant, up to null sets. We now return to the
left invariance under $N$. Take $a\in V_0$. By
Proposition~\ref{prop:finite-defect}, there are $r_j\in R\subset N$
converging to $a$ in the restricted adelic product away from $v_0$.
On the space obtained by forgetting the determinant-one variable at
$v_0$, the left rotations $L_{r_j}$ converge to $L_a$. This convergence
passes to bounded measurable functions in the required measure class:
restricted adelic convergence places all sufficiently late $r_j$ and
$a$ in the same integral compacts outside one fixed finite set; those
tail actions preserve measure. At the finitely many remaining local
factors, the smooth actions converge and their Radon--Nikodym densities
are uniformly bounded. Approximate a bounded function in $L^1$ by
continuous cylinder functions to pass invariance to the limit.
The determinant at $v_0$ is unchanged by every element of $S$, so it
causes no additional action on this reduced space.
Since $P$ is invariant under each $L_{r_j}$, it follows that
$P(L_a y)=P(y)$ almost everywhere for every $a\in V_0$.

Conditional equivariance and the left-rotation palette identity also give
$P(L_a y)=\phi(a)P(y)$. The image of $V_0$ is all of
$\mathcal F$, so $P(y)$ is the uniform probability almost everywhere.
For arbitrary $h$, the equality $Q_{hy}=h_*Q_y$ identifies the
$c_h$-marginal at $y$ with the $c_1$-marginal at $hy$. Nonsingularity
therefore makes every conditional coordinate marginal uniform.

\paragraph{Additive averaging.}
We have obtained all colors with equal conditional probability, but
the joint law is still conformal rather than translation invariant.
On the additive chart over a finite set $T$, replace $m_T$ in
Equation~\ref{eq:palette-disintegration} by additive Haar measure on
$J_T$, keeping $Q_y$ and the tail base $m^T$. The resulting
$\sigma$-finite joint measure is invariant under every rational
translation $T_b$ integral outside $T$: Haar measure is translation
invariant on $J_T$, the good integral translations preserve $m^T$, and
the conditional laws are equivariant.
Restrict to expanding Folner windows in $J_T$, divide by their Haar
volumes, and project to $\mathcal C$. Any weak limit is invariant under
all these translations. Every averaging measure already has uniform
coordinate marginals, because the conditional marginals were uniform.
Finally exhaust the places by such finite sets $T$ and take another
weak limit. Every $b\in J(k)$ is integral off a sufficiently large
$T$, so this limit has all the asserted invariances and marginals.
\end{proof}

\subsection{Fractional transformations and simultaneous color returns}

Fix the conjugacy-invariant subset $\mathcal S$ supplied by
Lemma~\ref{lem:finite-obstruction}. We shall relate the colors of $w$ and
$\tau_s w$ to returns under a quadratic additive translation.
For $w\in U(k)$ and $s\in k$, put
\begin{equation}\label{eq:palette-boost-data}
 t=\frac{s}{s-2i},\qquad C=tw,\qquad
 R_C=(1-C^*)(1-C)^{-1},\qquad z=wR_C.
\end{equation}
These expressions are defined for every $s$: $i\notin k$, and
\[
 d:=1-t\bar t=\frac{-4i^2}{s^2-4i^2}\ne0.
\]
The elements $C$ and $C^*$ commute, and $CC^*=t\bar t$ is central.
Division and $d\ne0$ imply that $1-C$ and $1-C^*$ are invertible.
It follows that $R_C$ is unitary and commutes with $w$. Hence $z$ is
unitary and commutes with $w$. Direct Cayley substitution gives
\begin{equation}\label{eq:palette-z-color}
 z=(w-\bar t)(1-tw)^{-1}
   =\frac{2i-s}{2i+s}\,\tau_s(w).
\end{equation}
The scalar factor is central and has norm one, so
$\phi(z)=\phi(\tau_s(w))$.

For such a $C$, the fractional transformation
\[
 D_C(u)=(u+C)(C^*u+1)^{-1}
\]
is induced by a form similitude: its matrix
$\left(\begin{smallmatrix}1&C\\C^*&1\end{smallmatrix}\right)$ has
multiplier $1-t\bar t$. The denominator is invertible for unitary $u$;
otherwise it vanishes, which would force $t\bar t=1$.
Writing $u=wa$ shows that
\begin{equation}\label{eq:palette-boost-commutes}
 w^{-1}D_C(u)=(a+t)(\bar t a+1)^{-1}
 \quad\text{commutes with }a=w^{-1}u.
\end{equation}
This calculation does not commute $u$ with $w$.

Starting from a variable $u\in U(k)$, define
\[
 v=D_{C^*}(uR_C),\qquad u'=R_C D_{-C}(v).
\]
Apply Equation~\ref{eq:palette-boost-commutes} first with unitary
parameter $w^{-1}$ and scalar $\bar t$. It gives
$[wv,wuR_C]=1$; conjugating by $w^{-1}$ gives $[vw,uz]=1$.
Applying the same equation with scalar $-t$ gives
$[w^{-1}v,z^{-1}u']=1$. Thus, writing
\[
 W=\phi(w),\quad Z=\phi(z),\quad
 X_1=\phi(u),\quad Y_1=\phi(v),\quad X_1'=\phi(u'),
\]
we have
\begin{equation}\label{eq:palette-two-commutators}
 [Y_1W,X_1Z]=1,\qquad [W^{-1}Y_1,Z^{-1}X_1']=1.
\end{equation}

The transformation $u\mapsto u'$ is an additive translation. Indeed its
fractional matrix is
\begin{align*}
 &\diag(R_C,1)
 \begin{pmatrix}1&-C\\-C^*&1\end{pmatrix}
 \begin{pmatrix}1&C^*\\C&1\end{pmatrix}
 \diag(1,R_C^{-1})\\
 &\hspace{35mm}=
 \begin{pmatrix}d+C-C^*&-(C-C^*)\\ C-C^*&d-(C-C^*)\end{pmatrix}.
\end{align*}
After dividing by $d$ and applying it to the column $(x+i,x-i)$,
the result is translation by $2i(C-C^*)/d$. Expanding the latter gives
\begin{equation}\label{eq:palette-quadratic-translation}
 b_w(s)=-(w+w^*)s-\frac{w-w^*}{2i}s^2\in J(k).
\end{equation}
The matrix calculation includes the point $\infty$ of the additive
chart. Neither $w$ nor the variable point is required to have a finite
Cayley coordinate.

Suppose that, for this fixed translation $T_{b_w(s)}$, every color
occurs as $X_1=X_1'$ at some rational point; different colors may use
different points. Put $B=Z^{-1}X_1$. Equations~\ref{eq:palette-two-commutators}
then imply
\[
 A:=W^{-1}Y_1\in C(B),\qquad
 D_0:=Z^{-1}Y_1WZ\in C(BZ^2).
\]
Since $WZ=ZW$, direct multiplication gives
\[
 A^{-1}(ZW^{-1})D_0=ZW.
\]
As $X_1$ ranges over $\mathcal F$, so does $B$, and this is exactly
Equation~\ref{eq:palette-obstruction}. Consequently
\begin{equation}\label{eq:palette-return-obstruction}
 \begin{gathered}
 \phi(w)\in\mathcal S,\quad\phi(\tau_s w)\notin\mathcal S\quad\Longrightarrow\\
 \text{not every color can recur under }T_{b_w(s)}.
 \end{gathered}
\end{equation}
When $s=0$, the translation is the identity and $Z=W$, so this implication
has no degenerate exception.

\subsection{Uniform finite witnesses for recurrence}

The translation-invariant law of Proposition~\ref{prop:palette-measure} will now
force simultaneous returns. Uniformity is essential: it lets us pass
the resulting restrictions from rational palettes to their closure.
A subset $D_0$ of the additive group $k$ is called \emph{thick} if it
contains a translate of every finite subset of $k$.

\begin{lemma}\label{lem:palette-uniform-recurrence}
Let $0<\delta\leq1$.
\begin{enumerate}[label=\textup{(\roman*)}]
\item There is an integer $r=r(\delta)\geq2$ such that, for every
probability-preserving action $T$ of $\bbZ^2$ and every event $E$ of
probability $\delta$, some $1\leq j\leq r$ satisfies
\[
 \mathbb P(E\cap T_{(j,j^2)}^{-1}E)>0.
\]
\item For every thick set $D_0\subset k$, there is a finite subset
$F_0\subset D_0$, depending only on $D_0,k,\delta$, such that for every
probability-preserving action $T$ of the discrete group $k^2$ and every
event $E$ of probability $\delta$, some $s\in F_0$ satisfies
\[
 \mathbb P(E\cap T_{(s,s^2)}^{-1}E)>0.
\]
\end{enumerate}
\end{lemma}

\begin{proof}
For an action of either abelian group, the spectral measure $\sigma$ of
$1_E$ is a positive measure on the compact character group, of total
mass $\delta$, with
\begin{equation}\label{eq:palette-spectral-atom}
 \sigma(\{1\})=\|\operatorname{proj}_{\rm inv}1_E\|_2^2
 \geq\delta^2.
\end{equation}
Its Fourier coefficient at $a$ is the return correlation
$\mathbb P(E\cap T_a^{-1}E)$, in particular a nonnegative real number.
The atom lower bound is preserved under weak limits: for the closed
singleton $\{1\}$ the limit mass is at least the upper limit of the
approximating masses.

For \textup{(ii)}, suppose no finite $F_0$ works. Enumerate $D_0$ and
choose a counterexample to each initial finite set. Compactness of the
spectral measures gives a limit $\sigma$ satisfying
Equation~\ref{eq:palette-spectral-atom} and having Fourier coefficient
zero at $(s,s^2)$ for every $s\in D_0$.
Take finite Folner sets $F_j\subset(k,+)$ and translates
$E_j=t_j+F_j\subset D_0$. We show that the average over $E_j$ of
\[
 f(s)=\chi_1(s)\chi_2(s^2)
\]
tends to zero for every nontrivial character pair $(\chi_1,\chi_2)$.
If $\chi_2=1$, this is ordinary averaging of a nontrivial linear
character. Such averaging is uniform in $t_j$, since a translate only
multiplies the average by a scalar of absolute value one.
If $\chi_2\ne1$, choose $M$ distinct shifts $a_1,\ldots,a_M\in k$.
Folner invariance allows replacement of the average of $f$ by the
average of $M^{-1}\sum_\ell f(s+a_\ell)$, with error tending to zero
independently of the translate. Cauchy--Schwarz, followed by expansion,
therefore gives
\begin{align*}
 \limsup_j\left|\frac1{|E_j|}\sum_{s\in E_j}f(s)\right|^2
 &\leq\frac1M+
 \frac1{M^2}\sum_{\ell\ne q}
 \limsup_j\left|\frac1{|E_j|}\sum_{s\in E_j}
 \chi_2(2(a_\ell-a_q)s)\right|\\
 &=\frac1M.
\end{align*}
All off-diagonal characters are nontrivial because multiplication by
$2(a_\ell-a_q)\ne0$ is onto $k$. Letting $M\to\infty$ proves the
claim. Dominated convergence against $\sigma$ now says that the average
of its coefficients at $(s,s^2)$ over $E_j$ tends to
$\sigma(\{1\})>0$, contradicting their vanishing.

For \textup{(i)}, failure of every bound $r$ similarly gives a measure
$\sigma$ on $\bbR^2/\bbZ^2$ satisfying
Equation~\ref{eq:palette-spectral-atom} and having coefficient zero at
every $(j,j^2)$ with $j\geq1$. Write a character phase as
$\exp(2\pi\mathrm i(\alpha j+\beta j^2))$, where here $\mathrm i$
is the usual complex imaginary unit. Choose finitely many pairs with
both $\alpha,\beta$ rational, including the trivial pair, so that the
remaining rational-pair mass is less than
$\varepsilon<\delta^2$. Choose a positive integer $M$ divisible by
their denominators, for example a sufficiently large factorial.
Along $j=M\ell$, all selected phases are $1$. Every pair with at least
one irrational coefficient has average zero by the degree-two form of
Weyl's theorem \cite{Weyl}. For the unselected rational pairs the
absolute contribution is bounded by $\varepsilon$; their individual
averages exist by periodicity. Dominated convergence makes the limiting
average at least $\delta^2-\varepsilon>0$, again a contradiction.
Increasing the resulting bound to at least $2$ is harmless.
\end{proof}

Apply the lemma to independent copies of $(\mathcal C,\mu)$ indexed by
$a\in\mathcal F$, and to the event
\[
 E=\{(c^{(a)})_{a\in\mathcal F}:c^{(a)}_1=a
       \text{ for every }a\in\mathcal F\}.
\]
Its probability is the fixed number
$\delta=|\mathcal F|^{-|\mathcal F|}>0$.
Writing $b_w(s)=sA_w+s^2B_w$ as in
Equation~\ref{eq:palette-quadratic-translation}, use the $k^2$-action
\[
 (a,b)\longmapsto T_{aA_w+bB_w}
\]
on this product space. Positive return of $E$ yields every recurrent
color in some palette. Each required two-coordinate pattern then
occurs rationally, so Equation~\ref{eq:palette-return-obstruction}
applies. For part~\textup{(i)}, and any $u_0\in k$, use instead the
$\bbZ^2$-action
\[
 (a,b)\longmapsto T_{a u_0A_w+b u_0^2B_w}.
\]
We have obtained two statements uniform in $w$:
\begin{enumerate}[label=\textup{(\alph*)}]
\item for every $u_0\in k$, some $s\in\{u_0,2u_0,\ldots,ru_0\}$
gives simultaneous recurrence of all colors;
\item for every fixed thick $D_0\subset k$, a fixed finite subset of
$D_0$ contains such an $s$ for every $w$.
\end{enumerate}
The second finite subset is not asserted to be uniform over thick sets.

\subsection{Additive invariance of membership in the finite subset}

Fix $c\in\mathcal C$ and a prefix $h\in\mathcal H$, and set
\[
 C_0=\{x\in k:c_{\tau_xh}\notin\mathcal S\}.
\]
If $x\notin C_0$, then $C_0-x$ cannot contain
$\{u_0,2u_0,\ldots,ru_0\}$ for any $u_0\in k$.
Indeed, match the finitely many relevant coordinates by a rational
palette and apply statement~\textup{(a)} to its rational point
$w=u(\tau_xhp)$. This contradicts
Equation~\ref{eq:palette-return-obstruction}.
Nor can $C_0-x$ be thick: apply statement~\textup{(b)} to that particular
thick set and match the finitely many witness coordinates by a rational
palette. The same contradiction results.

We record explicitly the combinatorial consequence:
\begin{equation}\label{eq:palette-density-alternative}
 C_0=k\quad\text{or}\quad C_0\text{ has zero upper Banach density}.
\end{equation}
For this purpose the upper Banach density of $C\subset k$ is the
supremum of $M(1_C)$ over translation-invariant means $M$ on bounded
functions on the discrete additive group $k$. Such means exist because
this group is abelian. If this density is zero, the density of $C$ along
every Folner sequence tends to zero: a nonzero upper limit would yield
an invariant mean with positive value by taking a weak limit of the
averaging functionals.

First, a run of $r$ consecutive points of $C_0$ in a direction $u_0$
extends to the adjacent point on either side. Otherwise that adjacent
point, lying outside $C_0$, would have the forbidden run immediately
ahead in direction $u_0$ or $-u_0$. Iteration fills the entire integer
progression.

Suppose now that $C_0$ has positive upper Banach density. To prove it
thick, fix any finite list $g_1,\ldots,g_d\in k$, and choose an invariant
mean $M$ with $M(1_{C_0})>\eta>0$. The finite multidimensional
Szemeredi theorem supplies a box size $N_0$ such that every subset of
$\{1,\ldots,N_0\}^d$ of density at least $\eta$ contains an affine
homothetic copy of $\{0,\ldots,r-1\}^d$, with positive integer dilation
$\ell\leq N_0$. This is the finite-box form of the density theorem
of \cite[Theorem~B]{FurstenbergKatznelson}.
Put $L_0=\operatorname{lcm}(1,\ldots,N_0)$ and average, with respect to
$M$, the density of the pullback of $C_0$ under
\[
 (m_1,\ldots,m_d)\longmapsto
 x+\sum_{j=1}^d m_jg_j/L_0.
\]
Translation invariance makes this average $M(1_{C_0})>\eta$, so one
translate has pullback density at least $\eta$. Injectivity of this
box map is not required. The resulting homothetic $r$-box lies in
$C_0$, with steps $\ell g_j/L_0$. Extend progressions successively along
each coordinate axis. This fills the whole integer grid with these
steps. Because $\ell$ divides $L_0$, that grid contains a translate of
$\{0,g_1,\ldots,g_d\}$. Thus $C_0$ is thick. If any $x$ were outside
$C_0$, then $C_0-x$ would also be thick, contrary to the preceding
restriction. This proves Equation~\ref{eq:palette-density-alternative}.

Now let the palette be random with law $\mu$. For a fixed Folner
sequence, Equation~\ref{eq:palette-density-alternative} makes the
densities of $C_0$ converge pointwise to $1_{\{C_0=k\}}$. Uniform
coordinate marginals and dominated convergence give
\[
 \mathbb P(C_0=k)=\frac{|\mathcal F\setminus\mathcal S|}{|\mathcal F|}.
\]
For each $s$, the event $C_0=k$ is contained in $s\in C_0$, and the
latter has exactly this same probability. Therefore
\begin{equation}\label{eq:palette-indicator-translation}
 1_{\mathcal S}(c_{\tau_sh})=1_{\mathcal S}(c_h)
 \quad\text{almost surely for every }s\in k,\ h\in\mathcal H.
\end{equation}
Countability makes these identities simultaneous for all $s,h$.
The palette rules also give, for every prefix $h$, invariance of this
indicator under left rotation by any element of $\ker\phi$ and under
conjugation by any element of $U(k)$, since $\mathcal S$ is
conjugacy-invariant. These are exact finite-coordinate identities, so
they hold throughout $\mathcal C$.

We have reached a concrete restriction on a common full-measure set of
palettes: membership of a color in $\mathcal S$ is unchanged under
scalar additive translations, kernel left rotations, and all unitary
conjugations, at every prefix. It remains to show that these
transformations already force invariance under all left colors.

\subsection{Generating left rotations by squares}

Let $\mathcal K$ be the subgroup of $\mathcal H$ generated by the three
classes of transformations just listed. Equation~\ref{eq:palette-indicator-translation}
and the simultaneous prefix identities show that
\begin{equation}\label{eq:palette-indicator-generated}
 1_{\mathcal S}(c_{ah})=1_{\mathcal S}(c_h)
 \qquad(a\in\mathcal K,\ h\in\mathcal H)
\end{equation}
on one set of full measure. This assertion is obtained by composing
finite-coordinate identities; it does not assume that $\mu$ is
invariant under the whole group $\mathcal K$.

Fix $\xi\in U(k)$. The commutative algebra $E=L(\xi)$ is a field, since
$D$ is division, and it is stable under $*$ because $\xi^*=\xi^{-1}$.
Let $E_0$ be its fixed field, so $E=LE_0$ and
$[E:E_0]=2$. Write $E^1$ for its norm-one group to $E_0$.
Choose $n_0\in E^1\cap\ker\phi$ whose finite Cayley coordinate
$x_0\in E_0$ generates $E_0$ over $k$. Such a choice is possible:
the rational norm-one torus is Zariski dense by its Cayley
parametrization, the power map on that torus is dominant, and powers
of an exponent killing $\mathcal F$ lie in $\ker\phi$. The requirement
that $x_0$ be finite and primitive is a nonempty Zariski-open condition.

Since there are only finitely many colors, choose $t_0\in k$ for which
the set
\[
 I=\{s\in k:\phi(\tau_sn_0)=\phi(\tau_{t_0}n_0)\}
\]
is infinite. For each $s\in I$, the word
\begin{equation}\label{eq:palette-affine-word}
 H_s=L_{n_0}^{-1}\tau_s^{-1}
 L_{(\tau_sn_0)(\tau_{t_0}n_0)^{-1}}
 \tau_{t_0}L_{n_0}
\end{equation}
belongs to $\mathcal K$. Both indicated left rotations have kernel
parameters. The word successively sends $1$ to $n_0$, to
$\tau_{t_0}n_0$, to $\tau_sn_0$, to $n_0$, and back to $1$.
It is represented by a fractional matrix over $E$; fixing $1$ in the
unitary chart means fixing infinity in the additive chart. Its latter
matrix is therefore upper triangular, say
\[
 \begin{pmatrix}\alpha&\gamma\\0&\beta\end{pmatrix},
 \qquad\alpha,\beta\in E^\times.
\]
On the commutative $E_0$-chart its affine multiplier is
\begin{equation}\label{eq:palette-affine-multiplier}
 a_s=\frac{\alpha}{\beta}
     =\frac{(x_0+t_0)^2-i^2}{(x_0+s)^2-i^2}\in E_0^\times.
\end{equation}
To check the multiplier, use
$q'(x)/q(x)=-2i/(x^2-i^2)$ in
Equation~\ref{eq:palette-affine-word}. The derivative at $u=1$ is
$((x_0+s)^2-i^2)/((x_0+t_0)^2-i^2)$; the affine multiplier at infinity
is its reciprocal. All denominators are nonzero because
$x_0+s\in E_0$ whereas $i\notin E_0$.

This commutative calculation produces transformations on the full
noncommutative chart, not merely on $E_0$. Conjugating $T_{s'}$ by the
upper triangular matrix gives the matrix
\[
 \begin{pmatrix}1&\alpha s'\beta^{-1}\\0&1\end{pmatrix}
 =\begin{pmatrix}1&s'a_s\\0&1\end{pmatrix}
 \qquad(s'\in k),
\]
since $s'$ is central. Thus $\mathcal K$ contains translations by every
$k$-multiple of $a_s$, and hence by their $k$-linear span.

That span is all of $E_0$. Indeed suppose a nonzero $k$-linear
functional on $E_0$ annihilated all the $a_s$ with $s\in I$.
Extend scalars to an algebraic closure and write this functional as a
linear combination of the distinct embeddings of $E_0$, which may be
extended while fixing $L$. Put $x_j$ for the corresponding distinct
conjugates of $x_0$. We obtain a linear combination, zero for infinitely
many $s$, of the rational functions
\begin{equation}\label{eq:palette-reciprocal-quadratics}
 \frac{(x_j+t_0)^2-i^2}{(x_j+s)^2-i^2}.
\end{equation}
It must vanish identically. These functions are linearly independent.
Their poles are $-x_j+i$ and $-x_j-i$. Poles from distinct denominators
can coincide only when the corresponding centers differ by $2i$ or
$-2i$. The resulting finite components form chains: each center has at
most one successor and predecessor, and characteristic zero excludes
cycles. An endpoint pole belongs to just one denominator, so its
coefficient in any linear relation is zero. Remove that denominator
and repeat along the chain. All numerators in
Equation~\ref{eq:palette-reciprocal-quadratics} are nonzero, so every
coefficient vanishes, the required contradiction. Consequently
$\mathcal K$ contains all $E_0$-translations.

Left rotation by $-1$ belongs to $\mathcal K$, because $-1$ is a central
kernel scalar. In the additive chart it is $x\mapsto i^2/x$, represented
by an antidiagonal matrix over $k$. It conjugates all upper
$E_0$-unipotents to all lower $E_0$-unipotents. These generate
$\SL_2(E_0)$, so $\mathcal K$ contains the corresponding fractional
transformations of the full chart. In particular it contains the
transformation whose matrix in the unitary chart is
$\diag(\xi,\xi^{-1})$: its conjugate by the Cayley change of basis is
\[
 \frac12
 \begin{pmatrix}
 \xi+\xi^{-1}&i(\xi-\xi^{-1})\\
 (\xi-\xi^{-1})/i&\xi+\xi^{-1}
 \end{pmatrix}\in\SL_2(E_0).
\]
Its action is $u\mapsto\xi u\xi$. Composing with the allowed conjugation
$u\mapsto\xi u\xi^{-1}$ gives $u\mapsto\xi^2u$. Hence
\begin{equation}\label{eq:palette-left-squares}
 L_{\xi^2}\in\mathcal K\qquad(\xi\in U(k)).
\end{equation}

\begin{proof}[Proof of Theorem~\ref{thm:no-simple-colors}]
All the preceding constructions were made under the assumption that
$\phi$ exists. Choose a palette satisfying the simultaneous identities
in Equation~\ref{eq:palette-indicator-generated}. Such palettes have
probability one by Proposition~\ref{prop:palette-measure} and the
recurrence argument. Equations~\ref{eq:palette-indicator-generated}
and~\ref{eq:palette-left-squares}, together with the exact left-rotation
palette rule, make membership in $\mathcal S$ invariant under left
multiplication by every square in $\mathcal F$, at all prefixes.
The subgroup generated by all squares is characteristic and nontrivial
in the nonabelian simple group $\mathcal F$: a group in which every
square is $1$ is abelian. Thus the squares generate $\mathcal F$.
Starting at the identity coordinate of the chosen palette, their left
multiplications run through every color. The membership indicator would
therefore have the same value on all of $\mathcal F$, contrary to
$\mathcal S$ being nonempty and proper. This contradiction proves the
theorem, subject to Lemma~\ref{lem:finite-obstruction}, whose proof follows.
\end{proof}
